\documentclass[11pt]{article}
\usepackage[a4paper,margin=28mm]{geometry}
\usepackage{amsmath,amsthm,amssymb}
\usepackage[hidelinks]{hyperref}
\newtheorem{theorem}{Theorem}
\title{Odd Cycle Transversal on subcubic graphs of large girth}
\author{Atlas\thanks{The proof in this note was found using ChatGPT.}}
\date{18 November 2026}
\begin{document}
\maketitle

In \textsc{Odd Cycle Transversal}, we ask whether deleting at most \(b\)
vertices of a graph \(G\) makes it bipartite. Write \(t\) for its
treewidth and \(B\) for the full input length, including any supplied
decomposition. A finite nonnegative exponent \(h\) is effectively
sublinear if there is a computable function \(T\) such that
\(h(u)\le u/j\) whenever \(j\ge1\) and \(u\ge T(j)\).

\hypertarget{odd-cycle-transversal-treewidth-w1}{}
\begin{theorem}
Assume \(\mathrm{FPT}\ne\mathrm{W}[1]\) and fix \(0<\eta<1\).
There is no deterministic algorithm for \textsc{Odd Cycle Transversal}
in time \(2^{h(t)}B^{\mathcal O(1)}\) with effectively sublinear \(h\),
even on connected simple graphs of maximum degree three and girth at
least \(n^{1-\eta}\), where \(n=|V(G)|\). This holds even when the
pathwidth equals \(t\) and an optimal path decomposition is supplied.
\end{theorem}

\begin{proof}
We reduce from \textsc{Clique} parameterized by the number \(k\) of
vertices sought, which is \(\mathrm{W}[1]\)-complete~\cite{df}.
Let the source graph have \(s\ge2\) vertices, and assume \(2\le k\le s\),
since the remaining cases are immediate. Use \(k\) blocks of
\(\lceil\log_2 s\rceil\) bits to name \(k\) vertices. Forbid each unused
name in each block by a clause, and forbid each equal or nonadjacent
pair of names in each pair of blocks by a clause. The resulting formula
is satisfiable exactly when the graph contains a \(k\)-clique.
It has \(m=k\lceil\log_2 s\rceil\) primary bits and length at most
\(A(k)s^c\) for a computable \(A\) and a constant \(c\).

\paragraph{A graph of width at most \(m+16\).}
We process the clauses one at a time to keep the path decomposition
small. Pad short clauses to length three by repeating literals. For
a clause of length \(r>3\), use the usual chain starting with
\((\ell_1\lor\ell_2\lor z_1)\), continuing with
\((\neg z_i\lor\ell_{i+2}\lor z_{i+1})\), and ending with
\((\neg z_{r-3}\lor\ell_{r-1}\lor\ell_r)\).
If every literal is false the chain forces successive links true and
violates its last clause. Otherwise, setting each link true exactly
when the prefix it represents is false satisfies the chain.

Introduce one common reference bit \(F\). Replace each ternary clause
\((a\lor b'\lor c')\) by the consecutive constraints
\(\operatorname{NAE}(a,b',y)\) and
\(\operatorname{NAE}(\neg y,c',F)\) with a fresh \(y\), where NAE
requires at least one value of each kind. When \(a=b'\) the first
constraint forces \(y=\neg a\), and the second fails precisely when
\(c'=F=a\). When \(a\ne b'\), some \(y\) always works. Complementing
all variable values preserves the constraints, so we may take \(F=0\).
Thus the pair represents the original clause.

Each constraint contributes a triangle with one vertex for each
literal position. Clause triangles for subcubic cut instances also
occur in Bonnet and Chang~\cite[Lemma~3.3]{bc}. Introduce a distinct
vertex for each occurrence of a variable. Link consecutive occurrences
of that variable, in processing order, by paths of length two. Link each
occurrence to its triangle vertex by a path of length two for a positive
literal and one edge for a negative literal. All path interiors are
private. The resulting simple graph has maximum degree three. The two
triangles of each pair are connected through \(y\), and the occurrence
links of \(F\) connect all pairs.

In any assignment of two colors to the vertices, every triangle has
at least one monochromatic edge. Equality between the number of such
edges in the whole graph and the number of triangles requires every
edge outside the triangles to be bichromatic. Consequently each variable
has one color on all its occurrences, and the triangle vertices carry the
literal values. A triangle contributes exactly one monochromatic edge
precisely when its constraint is satisfied. Conversely, any satisfying
assignment extends to a coloring attaining equality.
On a subcubic graph the minimum number of monochromatic edges equals
the minimum transversal size. Indeed, deleting an endpoint of each
monochromatic edge leaves a bipartite graph. In the other direction,
properly color the graph after deleting a transversal and reinsert its
vertices one at a time. At each step some color agrees with at most one
already colored neighbor. The threshold \(b\) is therefore the number of
triangles.

Make one bag per triangle. Carry the last occurrence of each variable
to its next occurrence. Include in the new bag its three occurrences,
its triangle, and all interiors of new sign and occurrence links.
Repeated occurrences within the bag can be linked there in order.
Each edge is covered, and each vertex occurs in an interval. Besides
the \(m\) primary variables, at most two successive chain variables,
the reference bit and the current \(y\) can be live. The new vertices
use at most twelve additional places. This gives width at most \(m+16\).

\paragraph{Large girth and a certificate of optimal width.}
Let \(q=m+40\). For an odd integer \(L\ge3\) to be chosen below,
replace every edge of the graph by a path of length \(L\). For fixed
endpoint colors such a path needs exactly one monochromatic edge if
the colors agree and zero otherwise. Subdivision thus preserves the
optimum and increases the girth to at least \(L\). At a bag containing
both endpoints, insert copies of that bag together with successive
pairs of subdivision vertices. This increases the width by at most two.

To force width exactly \(q\), add a bipartite component with a
\(K_{q+1}\) minor. Take \(q+1\) paths numbered \(0,\ldots,q\).
For each pair \(i<j\), join a designated vertex of path \(i\) to a
designated vertex of path \(j\) by an edge. Give every pair distinct
designated vertices and order them along each path by the lexicographic
order of the pairs. Prescribe black on the smaller-index endpoint of
each joining edge and white on the other. Separate consecutive designated
vertices on the same path by distance \(L\) or \(L+1\), choosing the parity
to extend the coloring. The component is bipartite and subcubic. Each
cycle traverses an entire such interval since the joining edges are a
matching. Contracting the paths exhibits the claimed clique minor.

Process pairs in lexicographic order, keeping the most recent vertex on
each unfinished path. Advance a path by introducing its next vertex and
forgetting the previous one. Keep both endpoints of the joining edge
in one bag. For a pair \((0,j)\), advance path zero before introducing
path \(j\). At most \(j+1\) vertices are then present (two for the first
pair). After the last pair of row zero retire path zero. At most \(q\)
paths remain and advancing one takes at most one extra place. This
decomposition has width at most \(q\), so the minor proves optimality.

Join an internal vertex of a subdivided edge of the first graph to the
first designated vertex of path zero by a bridge. Both endpoints had
degree at most two. Retain the first endpoint until the end of the first
decomposition, append a bag consisting of the bridge endpoints, and
start the second decomposition at the designated vertex. The first
decomposition still has width at most \(m+19\). The bridge creates no
cycle, so the budget is unchanged, and the combined graph has both widths
equal to \(q\).

Let \(H\) be the order before stretching, counting the \(q(q+1)\)
designated vertices of the second graph. Its final order \(n\) is at
most \(CHL\) for an absolute integer \(C\), by subcubicity and the
number of consecutive pairs. Fix an integer \(a\ge (1-\eta)/\eta\)
and choose the smallest odd \(L\ge\max\{3,(CH)^a\}\).
Then \(L^\eta\ge(CH)^{1-\eta}\), so the girth is at least
\(L\ge n^{1-\eta}\).
For fixed \(\eta\), the construction, including the listed bags and
their bit encoding, takes \(A'(k)s^{c'}\) time for computable \(A'\)
and constant \(c'\).

\paragraph{Running time.}
Suppose the claimed solver exists and let \(j=2(k+1)\).
If \(q\ge T(j)\), then \(h(q)\le q/j\le \log_2 s+21\).
Otherwise, since \(h(u)=o(u)\) and is finite, some integer \(D\)
satisfies \(h(u)\le u+D\) for all \(u\). Thus \(h(q)\le T(j)+D\)
in this second case. The two estimates bound the running time of
the resulting \textsc{Clique} algorithm by a computable function
of \(k\) times a fixed polynomial in \(s\), contradicting the hypothesis.
\end{proof}

\begin{thebibliography}{9}
\bibitem{bc} \'E. Bonnet and Y. Chang.
Max Independent Set Remains NP-hard when Excluding a Planar Induced
Minor (and More). \href{https://doi.org/10.48550/arXiv.2609.11285}{arXiv:2609.11285},
version 2, 2026.
\bibitem{df} R. G. Downey and M. R. Fellows.
Fixed-parameter tractability and completeness II: On completeness
for W[1]. \emph{Theoretical Computer Science} 141 (1995), 109--131.
\href{https://doi.org/10.1016/0304-3975(94)00097-3}{doi:10.1016/0304-3975(94)00097-3}.
\end{thebibliography}
\end{document}
